% : a natural-language task defines the goal, an isolated per-episode workspace carries observable and mutable state, a fixed set of domain-agnostic atomic tools constitutes the action space, a multi-turn reason--act--observe loop drives progress, and a terminal verifier scores the final workspace without prescribing operation paths.
% What remains is a five-element loop:
% (i)~a task instruction $q$ in natural language;
% (ii)~an isolated workspace $\mathcal{W}_0$, a per-trajectory file snapshot with no cross-sample sharing;
% (iii)~a fixed set of domain-agnostic atomic tools $\mathcal{A}=\{\texttt{read},\,\texttt{search},\,\texttt{edit},\,\texttt{exec},\,\texttt{finish}\}$ shared across all tasks;
% (iv)~a multi-turn interaction loop (reason--act--observe) until the model terminates or a budget $H$ is reached;
% (v)~a terminal state verifier $\mathcal{V}(\mathcal{W}_T)\in\{0,1\}$ that checks the final workspace without prescribing operation paths.
% Each episode is thus fully specified by a triple $\tau=(q,\,\mathcal{W}_0,\,\mathcal{V})$.
% (sparse terminal reward).
% The policy optimization objective is:
% \begin{equation}
% \label{eq:obj}
%     \max_{\theta}\;\mathbb{E}_{\tau\sim\pi_\theta}\!\big[\mathcal{V}(\mathcal{W}_T)\big].
% \end{equation}

% Full hyperparameters are summarized in Tab.~\ref{tab:hparam}.
% Rollouts are generated by a batched inference engine with the sampling temperature and top-$p$ listed in Tab.~\ref{tab:hparam}; 

% \begin{table}[t]
%     \centering
%     \small
%     \begin{tabular}{@{}ll@{}}
%     \toprule
%     \textbf{Hyperparameter} & \textbf{Value} \\
%     \midrule
%     Base model & Qwen3-8B \\
%     RL algorithm & GRPO \\
%     Global rollout batch & \textit{[e.g., 256]} \\
%     Mini-batch size & \textit{[e.g., 64]} \\
%     GRPO group size $G$ & \textit{[e.g., 8]} \\
%     Learning rate & \textit{[e.g., 1e-6]} \\
%     LR schedule & cosine, warmup \textit{[e.g., 0.03]} \\
%     Clip range $\epsilon$ & \textit{[e.g., 0.2]} \\
%     KL coefficient & \textit{[e.g., 0.0 / 0.001]} \\
%     Max context length & \textit{[e.g., 16384]} \\
%     Max turns $H$ & \textit{[e.g., 20]} \\
%     Max new tokens / turn & \textit{[e.g., 2048]} \\
%     Rollout temp / top-$p$ & \textit{[e.g., 1.0 / 1.0]} \\
%     Inference engine & \textit{[vLLM / SGLang]} \\
%     Sandbox concurrency $N$ & \textit{[e.g., 64]} \\
%     GPUs & 64 $\times$ \textit{[H100 / A100]} \\
%     Total wall-clock & \textit{[e.g., 36 h]} \\
%     \bottomrule
%     \end{tabular}
%     \caption{Training hyperparameters. Evaluation uses temperature $0.7$ over $3$ trials per instance.}
%     \label{tab:hparam}
% \end{table}



% \subsection{Analysis}
% \label{subsec:analysis}

% \paragraph{Sensitivity to $k_{\min}$.}
% The looping detector requires a minimum repeat length $k_{\min}$.
% Varying $k_{\min}\in\{8,16,32,64\}$ changes SR by less than \textbf{X.X} points (Tab.~\ref{tab:sensitivity}), indicating robustness to this hyperparameter.

% \paragraph{Qualitative comparison.}
% Fig.~\ref{fig:case} contrasts a GRPO-trained agent (post-collapse) with an AAM-trained agent on the same task.
% The GRPO agent exhibits a 12-turn looping pattern before truncation; the AAM agent completes the task in 5 turns with no repeated actions.

% \begin{table}[t]
% \centering
% \small
% \caption{Sensitivity to looping threshold $k_{\min}$ (Qwen3.5-9B, in-domain Avg.\ SR).}
% \label{tab:sensitivity}
% \begin{tabular}{ccccc}
% \toprule
% $k_{\min}$ & 8 & 16 & 32 & 64 \\
% \midrule
% SR (\%) & -- & -- & -- & -- \\
% \bottomrule
% \end{tabular}
% \end{table}

% \begin{figure}[t]
%   \centering
%   \includegraphics[width=\columnwidth]{./fig_case.pdf}
%   \caption{Qualitative comparison on a file-management task. Left: GRPO agent loops on redundant \texttt{ls} calls. Right: AAM agent executes a direct read--edit--verify sequence.}
%   \label{fig:case}
% \end{figure}